\section{Polynomial conditioning and an effective finite comparator (NS-81)}
\label{sec:ns81-effective-sampling}

\paragraph{Scope.}
The coefficient inverse in NS-78 makes the finite arithmetic comparator
in NS-80 effective. We prove an ordinary positive-Gram conditioning bound
and an explicitly sufficient finite sample cutoff. The cutoff is extremely
large; this is not a practical solver or a cofinal error estimate. None of
these bounds concerns a signed Weil operator or its spectral floor.

Let $H_N=\sum_{n=1}^N1/n$ in this section only, and let $G_N$ be the
ordinary Gram of $\sigma_1,\ldots,\sigma_N$. For a pure combination
$f=\sum c_n\sigma_n$, define the finite prefix form
\[
 P_N(c)=\eta^2+\sum_{m=1}^N\frac{f(1/m)^2+f(1/(m+1))^2}{m(m+1)},
 \qquad \eta=\sum_{n\leq N}c_n/n.
\]
It includes the sample at $1/(N+1)$. The same notation $P_N$ denotes
its matrix, and put
\[
 C_N=48H_NN^3(N+2).
\]

\begin{proposition}[A finite-prefix coefficient bound]
\label{prop:ns81-conditioning}
For every real coefficient vector,
\begin{equation}
 \|c\|_{\ell^2}^2\leq C_N P_N(c),\qquad
 \lambda_{\min}(G_N)\geq\frac1{16C_N}.
 \label{eq:ns81-conditioning}
\end{equation}
Writing $K_2=\|\sigma_1\|_2^2$, one also has
\[
 \lambda_{\max}(G_N)\leq K_2H_N,\qquad
 \operatorname{cond}_2(G_N)\leq768K_2H_N^2N^3(N+2)
                         =O(N^4\log^2(2N)).
\]
This is an ordinary spectral condition bound in the physical atom
coordinates, not a lower bound on approximation error.
\end{proposition}
\begin{proof}
Use $R=-f$ and $\alpha=0$ in NS-78; signs of the samples do not affect
$P_N$. Since $m\log(1+1/m)\geq1/2$, the exact inverse gives
\[
 \sum_{m=1}^N\frac{D_m^2}{m^2}
 \leq12\sum_{m=1}^N(z_m^2+z_{m+1}^2+\eta^2)
 \leq12N(N+2)P_N(c).
\]
Here each cell weight is at least $1/[N(N+1)]$ and $\eta^2\leq P_N$.
With $D_0=0$ and $\delta_m=D_m-D_{m-1}$,
\[
 \sum_{m=1}^N\frac{\delta_m^2}{m^2}
 \leq4\sum_{m=1}^N\frac{D_m^2}{m^2}
 \leq48N(N+2)P_N(c).
\]
M\"obius inversion and weighted Cauchy--Schwarz imply
\[
 |c_n|^2\leq\left(\sum_{d\mid n}d\right)
                  \sum_{d\mid n}\frac{\delta_d^2}{d}
 \leq nH_N\sum_{d\mid n}\frac{\delta_d^2}{d}.
\]
After summing and using $\sum_{k\leq y}k\leq y^2$ for $y\geq1$,
\[
 \sum_{n\leq N}|c_n|^2
 \leq H_N\sum_{d\leq N}\delta_d^2\sum_{k\leq N/d}k
 \leq H_NN^2\sum_{d\leq N}\frac{\delta_d^2}{d^2}
 \leq C_NP_N(c).
\]
The complete sampling form dominates its finite prefix, and NS-78
therefore gives $c^TG_Nc\geq P_N(c)/16$. This proves the ordinary
minimum-eigenvalue bound. The maximum is at most the trace,
$\operatorname{tr}G_N=K_2H_N$, by dilation and positivity. Division
of these two bounds proves the condition estimate.
\end{proof}

\begin{proposition}[Explicit complete tail and a finite comparator]
\label{prop:ns81-finite-comparator}
Let $Q^{<M}$ be the complete arithmetic sampling Gram of NS-80 truncated
to cells $1\leq m<M$, retaining the exterior term, where $M\geq N+1$.
Set
\[
 a_{N,M}=\tfrac12\log(2\pi M)+\frac{N+3}{6M},\qquad
 B_N(M)=\frac{2N}{M}\left(a_{N,M}^2+a_{N,M}+\tfrac12\right).
\]
Then the entire discarded sample tail satisfies
\begin{equation}
 0\preceq Q-Q^{<M}\preceq B_N(M)I
                  \preceq C_NB_N(M)Q^{<M}.
 \label{eq:ns81-tail-matrix}
\end{equation}
In particular the explicit integer choice
\begin{equation}
 M=2^{16}(N+1)^6
 \label{eq:ns81-explicit-cutoff}
\end{equation}
satisfies $C_NB_N(M)<1$, and hence
\[
 Q^{<M}\preceq Q\preceq2Q^{<M},\qquad
 \frac1{16}Q^{<M}\preceq G_N\preceq\frac{21}{8}Q^{<M}.
\]
The same constants survive any old-space Schur complement. For an
$N$-to-$2N$ block use~\eqref{eq:ns81-explicit-cutoff} with the total
atom count $2N$, not the old count $N$.
\end{proposition}
\begin{proof}
For $t\geq M\geq N$, the complete Bernoulli remainder in NS-78 gives
\[
 0\leq A(t/n)\leq\tfrac12\log(2\pi t)+\frac{N}{6M}
 \quad(1\leq n\leq N).
\]
On a cell $m\leq t\leq m+1$ with $m\geq M$, either endpoint's
atom value is bounded above by
$a_{N,M}+\tfrac12\log(t/M)$, since
$\log((m+1)/t)\leq1/M$. For a unit coefficient vector,
$\sum|c_n|\leq\sqrt N$, so both squared endpoint values are bounded
by $N(a_{N,M}+\tfrac12\log(t/M))^2$.
Integrating with $dt/t^2$ on every remaining cell gives $B_N(M)$,
including the complete tail to infinity. Finally
$Q^{<M}\succeq P_N\succeq I/C_N$, proving the last comparison.

For the displayed cutoff put $v=N+1\geq2$. Elementary bounds give
\[
 H_N\leq1+\log v\leq3v^{1/4},\qquad
 a_{N,M}\leq7+3\log v\leq13v^{1/4}.
\]
Here $\log v\leq2v^{1/4}$ follows by maximizing
$\log v/v^{1/4}$, whose maximum is $4/e<2$.
Also $\tfrac12\log(2\pi2^{16})<6.5$ and
$(N+3)/(6M)<1/2$, giving the bound for $a_{N,M}$.
Since $N^4(N+2)\leq v^5$, we obtain
\[
 C_NB_N(M)
 \leq\frac{96}{2^{16}}v^{-1}
       (3v^{1/4})\left(\frac{365}{2}v^{1/2}\right)
 =\frac{52560}{65536}v^{-1/4}<1.
\]
The full Gram comparison follows from NS-80. Taking old-coordinate
minima preserves the constants exactly as there.
\end{proof}

\paragraph{What is now finite, and what remains open.}
The finite comparator has projected condition ratio at most $42$.
The elementary NS-80 argument gives one-step capture at least
$168/1849$ of full available block gain and forty ideal steps capture
\[
 1-(41/43)^{80}>.9778.
\]
This is an effective finite-dimensional prescription; it does not make
that prescription fast. For a block with old size $256$, the total atom
count is $512$, and the sufficient cutoff exceeds $10^{21}$ cells.
No such enormous matrix is assembled in this update. The cutoff is a
conservative existence bound, not a necessary cutoff or a complexity
lower bound on every algorithm.

The scalar inequalities and positivity of the small finite-prefix matrices
are replayed with certified arithmetic. Those checks do not evaluate the
large finite comparator. Polynomial conditioning removes a numerical
resolution concern for these ordinary Grams, but leaves the missing
cofinal arithmetic decay completely unresolved. The original Weil floor,
G2 and RH remain open.

\paragraph{Finite validation.}
At $N=16,64,256$, the first $N$ cells and their final endpoint are
assembled exactly with Arb logarithms and integer factorial identities.
Their positive definiteness is supplied by the coefficient inverse.
A certified solve for $P_N^{-1}$ gives
$\operatorname{tr}(P_N^{-1})<C_N$ at each tested size, independently
confirming the displayed prefix floor; no midpoint solve is used.
The two 256/384-bit replays agree. At $N=256$ the analytic ordinary
Gram eigenvalue lower bound exceeds $4.91\cdot10^{-14}$, and the scalar
complete-tail ratio at~\eqref{eq:ns81-explicit-cutoff} is below $.01925$.
The latter is an evaluation of the proved tail bound, not an assembly
of the comparator's $18883333817345572864$ cells. Eighteen exact Maxima
checks validate the constants, primitives and rational step guarantees.
